% theory/pte/statements.tex -- a standalone LaTeX fragment. NOT wired into paper/main.tex.
%
% Assumes the preamble of paper/main.tex, the macros of theory/definitions.tex (\Orb, \Kmult),
% and theory/signatures/statements.tex (\Sig, lem:sigdata, thm:sigsep, cor:sigarea,
% thm:signonuniform, cor:siggrowth). It is meant to replace cor:siggrowth and the last sentence
% of thm:signonuniform's discussion.
% Citation keys: borweiningalls1994 (as in review/literature-pass/references-additions.bib);
% melzak1961, blp2003, cmsv2024, chen2025survey, wooley2019, crootmaoyip2026
% (entries in theory/pte/references-pte.bib).
% Full proofs: theory/pte/proof.md. Exact checks: witnesses.py, structure.py, growth.py,
% pencil_search.py, search_T3.sh in theory/pte.

\subsection{How many coefficients: the Prouhet--Tarry--Escott connection}\label{subsec:pte}

By Lemma~\ref{lem:sigdata} and Theorem~\ref{thm:sigsep}, two orbifolds in $\Sig$ with different
signatures share their first $L$ heat coefficients exactly when the cancelled mirror multiset
$Z=U^*\uplus(-V^*)$ is an \emph{$L$-configuration}. This means a nonempty multiset of nonzero
rationals, of even size and with no pair $\{z,-z\}$, such that
\[
  \sum_{z\in Z}z^{j}=0\quad(j\ \text{odd},\ 1\le j\le 2L-3),\qquad \sum_{z\in Z}z^{-1}=0 .
\]
Its size is $|Z|=|U^*|+|V^*|$, and its imbalance is $\iota(Z)=\#\{z>0\}-\#\{z<0\}=2(g'-g)$.

Let $N(k)$ be the least size of a non-trivial solution of the Prouhet--Tarry--Escott problem
of degree $k$, i.e.\ two distinct multisets of $n$ integers with equal power sums of exponents
$1,\dots,k$ \cite[\S1]{borweiningalls1994}. One has $k+1\le N(k)\le\tfrac12k(k+1)+1$
\cite[Props.~2, 3]{borweiningalls1994}. No bound $N(k)=o(k^2)$ is known, and this is listed
as an open problem in \cite[\S6]{borweiningalls1994}.

\begin{theorem}[Shape of a genus collision]\label{thm:ptedescartes}
Every $L$-configuration satisfies $|\iota(Z)|\le|Z|-2L$. Hence, if orbifolds in $\Sig$ of genera
$g\neq g'$ share their first $L$ heat coefficients, then
\[
  |U^*|+|V^*|\ \ge\ 2L+2|g-g'| ,
\]
and equality $|U^*|+|V^*|=2L+2$ forces $\{|U^*|,|V^*|\}=\{L,L+2\}$.
\end{theorem}

\begin{proof}[Sketch]
Write $Q(x)=\prod_{z\in Z}(x-z)$. Its odd elementary symmetric functions $e_k$ vanish for
$k\le2L-3$ and for $k=|Z|-1$. Since $|Z|$ is even, at most $(|Z|-2L)/2$ odd-degree coefficients
of $Q$ are nonzero. All roots are real and nonzero, so Descartes' rule of signs is an equality
for $Q(x)$ and $Q(-x)$. The imbalance is then the signed number of odd gaps between consecutive
nonzero coefficients, and each odd-degree coefficient creates at most two of them.
\end{proof}

\begin{proposition}[Symmetric constructions do not change the genus]\label{prop:ptebalanced}
Let $A$ be an odd ideal symmetric solution of size $2L-1$, i.e.\ $\sum_{a\in A}a^j=0$ for odd
$j\le2L-3$ \cite[p.~8]{borweiningalls1994}. Then
$\#\{a>0\}-\#\{a<0\}=\operatorname{sgn}\sum_{a\in A}a^{-1}$. Consequently, for two such sets
$A,B$, every $L$-configuration $A\uplus\lambda B$ ($\lambda\in\mathbb Q^\times$) has
$\iota=0$.
\end{proposition}

\begin{theorem}[Growth]\label{thm:ptegrowth}
Let $f(A)=\max\{\Kmult(\Orb;\Sig):\operatorname{Area}(\Orb)\le A\}$.
\begin{enumerate}[label=(\alph*)]
\item For $A\ge8\pi$,
  \[
    \Bigl\lfloor\sqrt{\tfrac13\bigl(\tfrac{A}{2\pi}-1\bigr)}\Bigr\rfloor+2\ \le\ f(A)\ \le\ \Bigl\lfloor\frac{A}{\pi}\Bigr\rfloor+4 .
  \]
\item If $f(A)\ge L+1$, then $N(2L-2)\le2\lfloor A/\pi\rfloor+8$.
\item For $0<\alpha\le1$: $f(A)\ge cA^{\alpha}$ for all large $A$ and some $c>0$ if and only if
  $N(k)\le Ck^{1/\alpha}$ for all $k$ and some $C$. In particular $f(A)=\Theta(A)$ if and only
  if $N(k)=O(k)$.
\end{enumerate}
The lower bound in (a) also holds, up to the constant, for the number of coefficients needed
to determine the genus alone, and for the number needed to determine the cone count alone
within genus $0$.
\end{theorem}

\begin{proof}[Sketch]
(a) Upper bound: Corollary~\ref{cor:sigarea}. Lower bound: by pigeonhole there are distinct
$n$-multisets $X,Y$ of positive integers with $n\le(L-1)^2+1$ and equal power sums of every
odd exponent $j\le2L-3$. Put
\[
  U=X\uplus2Y\uplus2Y,\qquad V=Y\uplus2X\uplus2X .
\]
Then
$\sum_Uu^j-\sum_Vv^j=(1-2^{j+1})\bigl(\sum_Xx^j-\sum_Yy^j\bigr)$, which vanishes for odd
$j\le2L-3$ and for $j=-1$. So $(0;U)$ and $(0;V)$, with any 1s removed, are distinct hyperbolic
orbifolds of area $<2\pi(3(L-1)^2+1)$ sharing $L$ coefficients.

(b) A configuration $Z$ of size $T$ gives the PTE solution $[Z]=_{2L-2}[-Z]$ of size $T$, and
$T\le2\lfloor A/\pi\rfloor+8$ by the proof of Corollary~\ref{cor:sigarea}.

(c) Combine (a) and (b), using that $N$ is nondecreasing. For the genus, shift a PTE solution
$X,Y$ of degree $2L-3$ into $(X+c)\uplus(-(Y+c))$. For the cone count, translate it so that it
contains $1$ and apply the doubling above.
\end{proof}

\begin{remark}[What the exponent means]\label{rem:pteexponent}
All known bounds on $N(k)$ are quadratic \cite[Prop.~3 and p.~7]{borweiningalls1994}, and with any
quadratic bound Theorem~\ref{thm:ptegrowth}(c) gives exactly the exponent $\tfrac12$ of (a). Any exponent $\alpha>\tfrac12$ would give
$N(k)=O(k^{1/\alpha})$, which would settle in a strong form the open problem
$N(k)=o(k^2)$ \cite[\S6, Q3]{borweiningalls1994}. Linear growth of $f$ is equivalent to
$N(k)=O(k)$, which is implied by the conjecture $N(k)=k+1$. Ideal solutions are known only for
$k\le9$ and $k=11$ \cite{cmsv2024,crootmaoyip2026}.
\end{remark}

\begin{example}[Small pairs]\label{ex:ptepairs}
The following pairs share exactly $L$ heat coefficients (exact computation with the cone
coefficients):
\begin{itemize}
\item $L=3$, same genus and cone count: $(0;3,10,15,30)$ and $(0;4,5,21,28)$, of area
  $2\pi\cdot\frac{22}{15}$. Previously the smallest area known to share three coefficients was
  $2\pi\cdot\frac{14}5$, for $(1;15,15,15)$ and $(0;3,3,5,7,7,21)$.
\item $L=3$, genus $0$ with $7$ and $8$ cone points: $(0;4,4,5,5,6,12,12)$ and
  $(0;2,2,2,3,10,10,10,10)$, obtained from $[1,5,5]=[2,3,6]$ \cite{chen2025survey} by the doubling
  above. This replaces $103$ vs $104$.
\item $L=4,5,6,7$: genus-$0$ pairs with equal cone counts and area
  $<2\pi\cdot5,\,7,\,10,\,18$. They come from pairs of odd ideal symmetric solutions of sizes $7$
  and $9$ \cite{borweiningalls1994,blp2003}, from equal sums of odd powers
  \cite{chen2025survey}, and from the size-$12$ ideal solution \cite{cmsv2024}. The
  Prouhet--Thue--Morse pairs of Theorem~\ref{thm:signonuniform} need area up to
  $2\pi(4^{L-1}-1)$.
\item Genus $1$ versus genus $0$ sharing $L=4,5,6,7$ coefficients: configurations of size
  $16,20,26,40$, against $2^{2L-1}=128,512,2048,8192$ for Thue--Morse.
\end{itemize}
\end{example}

\begin{remark}[The first open case]\label{rem:ptet3}
By Theorem~\ref{thm:ptedescartes}, a genus collision sharing three coefficients with
$|U^*|+|V^*|=8$ must have shape $(3,5)$. An exhaustive exact search excludes all such collisions
whose five-element side, made primitive, has entries $\le220$. Real solutions of this shape
exist. So whether the least size is $8$ or $10$ remains open.
\end{remark}

% Exact extent. thm:ptedescartes, prop:ptebalanced, thm:ptegrowth are proved for all L given the
% heat input of lem:sigdata (theory/pte/proof.md sections 2-4); ex:ptepairs is exact computation
% (witnesses.py). The bound N(k) <= k(k+1)/2+1 is cited (borweiningalls1994 Prop. 3) and also
% reproved by the same pigeonhole in proof.md Lemma 1.5.
